\clearpage
\sscp{Innovations in the lexicon}{07-04-02}
Lexically specific innovations that distinguish the subgroups
of central Maluku are summarised below.
The few exceptions are attributable to known contact.
There is some overlap with the previous discussion about the
lexically specific patterns in the historical phonology.

\begin{enumerate}
	\item	*j > {\0} in \tsc{Sula-Buru} in *ŋajan ‘name’,
				and *huaji ‘younger sibling (same sex)’,
				but is preserved as *j > **l in \tsc{\ili{Ambon}-Seram} languages in
				the same words (\trf{07-12}).
	\item	PMP *kahiw ‘wood,
				tree’ > **kau in \tsc{Sula-Buru} languages; > **kai in \tsc{\ili{Ambon}-Seram} and \tsc{Seram-Tanimbar-Bomberai}.
				(See \trf{07-13}.) \ili{Kayeli} follows the \tsc{Sula-Buru} pattern due to contact.
				Since \tsc{Sula-Buru} languages \ili{Sanana}
				and \ili{Mangole} normally preserve *y = \it{y},
				it appears they did not derive their words for
				‘tree’ from purported PCEMP **kayu.
	\item	\tsc{Sula-Buru}
				and \tsc{Seram-Tanimbar-Bomberai} languages take their term for
				‘boat, (outrigger) canoe’ from a *w-initial form,
				as from PCEMP **waŋka.
				\tsc{\ili{Ambon}-Seram} languages take their term from a *b-initial
				form, as from PMP *baŋkaq ‘boat’.
				See \trf{07-13}.
				\ili{Kayeli} follows the \tsc{Sula-Buru} pattern due to contact.
	\item	All \tsc{Sula-Buru} languages take their word for ‘eye,
				source’ from Proto-Sula-\ili{Buru} **dama.
				Other  languages of Maluku take their word from PMP *mata. (See \trf{07-13}.)
				\ili{Kayeli} follows the \tsc{Sula-Buru} pattern due to contact.
	\item	All \tsc{Sula-Buru} languages take their word for ‘banana’ from **biak.
				\tsc{\ili{Ambon}-Seram} languages for which we have data take their
				word for ‘banana’ from PMP *punti (seven languages),
				**kula (six languages), and **telewa (four languages).
				\tsc{Seram-Tanimbar-Bomberai} languages take their word for ‘banana’
				from *punti or {\#}muku.
				No \tsc{\ili{Ambon}-Seram} or \tsc{Seram-Tanimbar-Bomberai} languages
				for which we have data take their word for ‘banana’ from **biak.
				No \tsc{Sula-Buru} or \tsc{\ili{Ambon}-Seram} languages for which we
				have data take their word for ‘banana’ from Papuan {\#}muku.
				See the broader discussion in \srf{02-03-01-07},
				local data in \trf{07-13},
				and banana-term data from 311 languages in appendix \ref{ch:BanTer}.
	\item	\tsc{Sula-Buru} numerals ‘one’,
				‘four’, and ‘six’ are shaped CVV for PMP numerals shaped *(C)VCV(C).
				\tsc{\ili{Ambon}-Seram} numerals ‘one’, ‘four’,
				and ‘six’ are shaped (C)VCV.
				\tsc{Seram-Tanimbar-Bomberai} numeral ‘one’ is shaped (V)CV,
				‘four’ is shaped CVC(V)
				and six is shaped either VCVC or CVVC.
				See \trf{07-14} and \trf{07-15}.
	\item	\tsc{Nuclear Sula-Buru} languages take their word for ‘ten’
				from PMP *puluq; \tsc{\ili{Ambon}-Seram} languages take their word
				for ‘ten’ from **butu.
				Most \tsc{Seram-Tanimbar-Bomberai} languages take their word
				for ‘ten’ from **butu,
				with the exception of members of \tsc{East Bomberai} (and possibly
				\ili{Kur}) which take their word for ‘ten’ from *sa-ŋa-puluq. See \trf{07-15}.
	\item	\tsc{\ili{Ambon}-Seram} languages take their word for ‘rat,
				mouse’ from a form *malabaw.
				No \tsc{Sula-Buru} or \tsc{Seram-Tanimbar-Bomberai} languages
				for which we have data derive from that form. See \trf{07-10}.
	\item	\tsc{PAn} *um-aRi,
				*um-ai ‘come’ > \it{mahi} in \ili{Buru}
				and \ili{Lisela} (*R > \it{h} is regular); but > \it{mai {\tl} mae}
				in \tsc{\ili{Ambon}-Seram} languages that normally preserve *R.
\end{enumerate}

In the social structure
and sibling terminology, \tsc{Sula-Buru} languages have four
siblings terms, including two cross-sex sibling terms: *bətaw ‘sister (male speaking)’
and *ñaRa ‘brother (female speaking)’.
\tsc{\ili{Ambon}-Seram} languages have three sibling terms,
with only one term *bətaw innovating the meaning ‘opposite sex
sibling’ (or **leku ‘opposite sex sibling’ innovated in the \tsc{\ili{Piru} Bay} branch).
See data and discussion in \srf{07-03}.
Note that four sibling terms are also found in other subgroups
as well, and some languages in other subgroups have also reduced to three
siblings terms (some using reflexes of *ñaRa instead of *bətaw
to mean ‘opposite sex sibling’.) \tsc{\ili{Ambon}-Seram} seems to be
the only one innovating as a whole subgroup on this issue.

\trf{07-13} shows patterned lexical differences between \tsc{Sula-Buru}
and \tsc{\ili{Ambon}-Seram} for ‘tree,
wood’, ‘boat, (outrigger) canoe’, ‘eye,
source’, and ‘banana’.

\begin{table}[p]\centering\tcp{\ili{Central} Maluku lexical differences}{07-13}
\begin{adjustbox}{totalheight=\textheight}
	\begin{threeparttable} %\st{2pt}
	\begin{tabular}{llllll}\lsptoprule
*gloss	&	tree, wood	&	boat, canoe	&	eye, source	&	 \mc{2}{l}{banana}\\
PMP	&	\hr*kahiw	&	\hr*baŋkaq	&	\hr*mata	&	\mc{2}{l}{\hr*punti, etc.}\\ \midrule
\tsc{Sula-Buru}	&	**kau	&	**waŋka	&	**dama	&	\mc{2}{l}{**biak}\\ \midrule
\ili{Mangole}	&	{\hr\hr}kau	&	\hr\cg{(lotu)}	&	{\hr\hr}lama	&	\mc{2}{l}{{\hr\hr}fia}\\
\ili{Sanana}	&	{\hr\hr}kau	&		&	{\hr\hr}hama	&	\mc{2}{l}{{\hr\hr}fia}\\
\ili{Lisela}	&	{\hr\hr}kau	&	{\hr\hr}waga	&	{\hr\hr}rama-n	&	\mc{2}{l}{{\hr\hr}fia-t}\\
\ili{Buru}	&	{\hr\hr}kau	&	{\hr\hr}waga	&	{\hr\hr}rama-n	&	\mc{2}{l}{{\hr\hr}fie-t\su{†}}\\
\ili{Ambelau}	&	{\hr\hr}ʔau-a	&	{\hr\hr}waʔa, βaa	&	{\hr\hr}lama-ni	&	\mc{2}{l}{{\hr\hr}bia-e}\\\midrule
\tsc{\ili{Ambon}-Seram}	&	**kai	&	**baŋkaq	&	**mata	&	\mc{2}{l}{\hr\hr(various)}\\ \midrule
\ili{Kayeli}	&	{\hr\hr}au (<SB)	&	{\hr\hr}waa (<SB)	&	{\hr\hr}lama-ni (<SB)	&	\mc{2}{l}{{\hr\hr}mpulu}\\
\ili{Manipa}	&	{\hr\hr}he-n	&		&	{\hr\hr}maka	&	\mc{2}{l}{}\\
\ili{Kelang}	&	{\hr\hr}hai	&		&		&	\mc{2}{l}{}\\
\ili{Luhu}	&	{\hr\hr}ai	&		&	{\hr\hr}mata	&	\mc{2}{l}{{\hr\hr}huli-i}\\
\ili{Piru}	&	{\hr\hr}ai	&	{\hr\hr}haka	&	{\hr\hr}mata	&	\mc{2}{l}{{\hr\hr}tema}\\
\ili{Larike}	&	{\hr\hr}ei	&	{\hr\hr}haka	&	{\hr\hr}ma'tesi	&	\mc{2}{l}{{\hr\hr}ure}\\
\ili{Asilulu}	&	{\hr\hr}ai	&	{\hr\hr}haka	&	{\hr\hr}mata-n	&	\mc{2}{l}{{\hr\hr}kula}\\
Kaitetu	&	{\hr\hr}ai	&	{\hr\hr}haka	&	{\hr\hr}mata-ɲ	&	\mc{2}{l}{{\hr\hr}kula}\\
\ili{Paulohi}	&	{\hr\hr}ai	&	\hr\cg{(sapan-e)}	&	{\hr\hr}mata	&	\mc{2}{l}{{\hr\hr}fuki}\\
\ili{Kaibobo}	&	{\hr\hr}ai	&	\hr\cg{(tala)}	&	{\hr\hr}mata	&	\mc{2}{l}{{\hr\hr}tema}\\
\ili{Kamarian}	&	{\hr\hr}ai	&	\hr\cg{(tala)}	&	{\hr\hr}mata	&	\mc{2}{l}{{\hr\hr}uki}\\
\ili{Haruku}	&	{\hr\hr}ai	&	\hr\cg{(tala)}	&	{\hr\hr}mata	&	\mc{2}{l}{{\hr\hr}kura}\\
\ili{Saparua}	&	{\hr\hr}ai	&	\hr\cg{(talal)}	&	{\hr\hr}mata	&	\mc{2}{l}{{\hr\hr}kula-l}\\
\ili{Nusalaut}	&	{\hr\hr}ai	&	\hr\cg{(talal)}	&	{\hr\hr}maa	&	\mc{2}{l}{{\hr\hr}telewa-l}\\
\ili{Sepa}	&	{\hr\hr}yai	&	{\hr\hr}haka	&	{\hr\hr}mata-n	&	\mc{2}{l}{{\hr\hr}telewa}\\
\ili{Tamilou}	&	{\hr\hr}yai	&	{\hr\hr}haka	&	{\hr\hr}mata	&	\mc{2}{l}{{\hr\hr}telewa}\\
\ili{Tehoru}	&	{\hr\hr}yai	&	{\hr\hr}haka	&	{\hr\hr}mata	&	\mc{2}{l}{{\hr\hr}telewa}\\
N. \ili{Wemale}	&	{\hr\hr}yai	&	\hr\cg{(kalutu)}	&	{\hr\hr}mata	&	\mc{2}{l}{{\hr\hr}huri}\\
N. \ili{Alune}	&	{\hr\hr}ai	&	\hr\cg{(soa buini)}	&	{\hr\hr}mata	&	\mc{2}{l}{{\hr\hr}uri, tema}\\
S. \ili{Nuaulu}	&	{\hr\hr}ai	&	{\hr\hr}haka	&	{\hr\hr}mata-i	&	\mc{2}{l}{{\hr\hr}uri}\\
\ili{Yalahatan}	&	{\hr\hr}ai	&	{\hr\hr}haka	&	{\hr\hr}mata-i	&	\mc{2}{l}{{\hr\hr}uta}\\
\ili{Manusela}	&	{\hr\hr}ai	&	\hr\cg{(palutu-a)}	&	{\hr\hr}mata	&	\mc{2}{l}{{\hr\hr}teloa}\\
\ili{Liana-Seti}	&	{\hr\hr}ai-a	&	{\hr\hr}waʔa \hr\cg{(loan?)}	&	{\hr\hr}mata	&	\mc{2}{l}{}\\ \midrule
\ili{Banda}	&	\hr\cg{(uno)}	&	\hr\cg{(farau, bot)}	&	{\hr\hr}mata-	&	\mc{2}{l}{{\hr\hr}kula}\\ \midrule
STB	&	**kai	&	**waŋka	&	**mata	&	**fundi, & **muku\\\midrule
\ili{Bobot}	&	{\hr\hr}ai	&	{\hr\hr}waa-n	&	{\hr\hr}mata-n	&	&\\
\ili{Masiwang}	&	{\hr\hr}ai	&		&	{\hr\hr}mata	&	\hr\cg{(tan)}&\\
\ili{Geser}	&	{\hr\hr}kai	&	{\hr\hr}aŋ	&	{\hr\hr}mata	&	{\hr\hr}fudi&\\
\ili{Gorom}	&	{\hr\hr}ai	&	{\hr\hr}aŋ	&	{\hr\hr}mata	&	{\hr\hr}hudi&\\
\ili{Watubela}	&	{\hr\hr}kai	&		&	{\hr\hr}mata	&	{\hr\hr}fudi&\\
\ili{Kowiai}	&	{\hr\hr}ai	&	\hr\cg{(biyer)}	&	{\hr\hr}mataɸut	&	{\hr\hr}ɸun&\\
\ili{Kur}	&	{\hr\hr}kai	&		&	{\hr\hr}matin	&	&\hr‹muk›\\
\ili{Yamdena}	&	{\hr\hr}ka-	&	\hr\cg{(sori)}	&	{\hr\hr}mata-n	&	{\hr\hr}funri&\\
\ili{Fordata}	&	{\hr\hr}aa	&	\hr\cg{(raa)}	&	{\hr\hr}mata-n	&	&{\hr\hr}muʔu\\
\ili{Kei}	&	{\hr\hr}ʔay	&	\hr\cg{(raw)}	&	{\hr\hr}mat-n	&	&{\hr\hr}muk\\
\ili{Uruangnirin}	&	{\hr\hr}ain	&	\hr\cg{(rau)}	&	{\hr\hr}mata puan	&	&{\hr\hr}maŋgo\\
\ili{Sekar}	&	{\hr\hr}yain	&	\hr\cg{(rai)}	&	{\hr\hr}bati-n	&	{\hr\hr}fudi&\\
\ili{Arguni}	&	{\hr\hr}ae, á	&	\hr\cg{(fúàg)}	&	{\hr\hr}tapútu	&	{\hr\hr}fúd&\\
\ili{Irarutu}	&	{\hr\hr}e	&	\hr\cg{(ʤu)}	&	{\hr\hr}mtie	&	{\hr\hr}fud&\\
		\lspbottomrule
	\end{tabular}
		\begin{tablenotes}\tnsize
			\item[†] {\ili{Buru} Wae Sama dialect \it{fie-t} ‘banana’.}
		\end{tablenotes}
	\end{threeparttable}
\end{adjustbox}
\end{table}
